\documentclass{beamer}
\usepackage[utf8]{inputenc}
\usepackage{graphicx}

\newtheorem{definicion}{Definición}
\newtheorem{ejemplo}{Ejemplo}

%%%%%%%%%%%%%%%%%%%%%%%%%%%%%%%%%%%%%%%%%%%%%%%%%%%%%%%%%%%%%%%%%%%%%%%%%%%%%%%
\title[Presentación con Beamer]{Presentación sobre $pi$ usando BEAMER}
\author[María Jesús]{María Jesús Álvarez Rodríguez}
\date[23-04-2014]{23 de abril de 2014}
%%%%%%%%%%%%%%%%%%%%%%%%%%%%%%%%%%%%%%%%%%%%%%%%%%%%%%%%%%%%%%%%%%%%%%%%%%%%%%%

%\usetheme{Madrid}
%\usetheme{Antibes}
%\usetheme{tree}
%\usetheme{classic}

%%%%%%%%%%%%%%%%%%%%%%%%%%%%%%%%%%%%%%%%%%%%%%%%%%%%%%%%%%%%%%%%%%%%%%%%%%%%%%%
\begin{document}
  
%++++++++++++++++++++++++++++++++++++++++++++++++++++++++++++++++++++++++++++++  
\begin{frame}

  \includegraphics[width=0.15\textwidth]{img/ullesc}
  \hspace*{7.0cm}
  \includegraphics[width=0.16\textwidth]{img/fmatesc}
  \titlepage

  \begin{small}
    \begin{center}
     Facultad de Matemáticas \\
     Universidad de La Laguna
    \end{center}
  \end{small}

\end{frame}
%++++++++++++++++++++++++++++++++++++++++++++++++++++++++++++++++++++++++++++++  

%++++++++++++++++++++++++++++++++++++++++++++++++++++++++++++++++++++++++++++++  
\begin{frame}
  \frametitle{Índice}  
  \tableofcontents[pausesections]
\end{frame}
%++++++++++++++++++++++++++++++++++++++++++++++++++++++++++++++++++++++++++++++  


\section{Primera Sección}


%++++++++++++++++++++++++++++++++++++++++++++++++++++++++++++++++++++++++++++++  
\begin{frame}

\frametitle{Primera Sección}

\begin{definicion}
El número PI equivale a la constante que relaciona el perímetro o longitud de una 
circunferencia con su diámetro. Se trata de un valor con un infinito número de decimales, cuya secuencia comienza de 
la siguiente manera:3,1415926535897932384626433832795028841

\end{definicion}

\end{frame}
%++++++++++++++++++++++++++++++++++++++++++++++++++++++++++++++++++++++++++++++  

\section{Segunda Sección}

%++++++++++++++++++++++++++++++++++++++++++++++++++++++++++++++++++++++++++++++  
\begin{frame}

\frametitle{Segunda Sección}

\begin{block}{Historia}
 
El número pi es la constante que relaciona el perímetro de una circunferencia,
con la longitud de su diámetro p = L/D. Este no es un número exacto, sino que es de los llamados números irracionales 
que tiene infinitas cifras decimales sin repetición de períodos. Ya en la antigüedad, se insinuó que todos los círculos 
conservaban una estrecha dependencia entre el contorno y su radio, pero tan sólo desde el siglo XVII la correlación se 
convirtió en un dígito y fue identificado con el nombre "Pi". A lo largo de la historia, a este ilustre guarismo se le han asignado diversas cantidades.
Sin embargo, fue en Grecia donde la correspondencia entre el radio y la longitud de una circunferencia comenzó a consolidarse 
como uno de los más insignes enigmas a resolver. Un coetáneo de Sócrates, Antiphon, inscribió en el círculo un cuadrado, 
luego un octógono e ideó multiplicar la cantidad de lados hasta el momento en que el polígono obtenido se ajustara casi con
la circunferencia.

\end{block}

\end{frame}

%++++++++++++++++++++++++++++++++++++++++++++++++++++++++++++++++++++++++++++++   


\section{Tercera Sección}


%++++++++++++++++++++++++++++++++++++++++++++++++++++++++++++++++++++++++++++++  
\begin{frame}

\frametitle{Tercera Sección}

\begin{definicion}
\[ S_n=a_1+\cdots + a_n = \sum_{i=1}^n a_i \]

\[\int_{x=0}^{\infty} x\,\text{e}^{-x^2}\text{d}x=\frac{1}{2},\quad\text{e}^{i\pi}+1=0\]

\[ \sqrt 2 = 1+\dfrac{1}{2+\dfrac{1}{2+\dfrac{1}{2+\dfrac{1}{\ddots}}}} \]

\[x=\frac{a_2 x^2 + a_1 x + a_0}{1+2z^3}, \quad x+y^{2n+2}=\sqrt{b^2-4ac}\]

\[ S_n=a_1+\cdots + a_n = \sum_{i=1}^n a_i \]

\end{definicion}

\end{frame}
%++++++++++++++++++++++++++++++++++++++++++++++++++++++++++++++++++++++++++++++ 

\section{Bibliografía}
%++++++++++++++++++++++++++++++++++++++++++++++++++++++++++++++++++++++++++++++  
\begin{frame}
  \frametitle{Bibliografía}

  \begin{thebibliography}{10}

    \beamertemplatebookbibitems 
    {\small $ http://www.saberia.com/2010/03/que-es-el-numero-pi/$}

    \beamertemplatebookbibitems
    {\small $http://www.monografias.com/trabajos57/calculo-pi/calculo-pi.shtm$}

  \end{thebibliography}
\end{frame}

%++++++++++++++++++++++++++++++++++++++++++++++++++++++++++++++++++++++++++++++  
\end{document}
